\PassOptionsToPackage{hyphens}{url}
\documentclass[11pt]{amsart}
\usepackage[margin=1.1in]{geometry}
\usepackage{amssymb,amsmath,booktabs,hyperref}
\newtheorem{theorem}{Theorem}
\theoremstyle{remark}
\newtheorem{remark}[theorem]{Remark}
\title{A remark on OpenAI's Smith--Toda complex $V(4)$:\\ from $p=1009$ to every prime $p\ge11$}
\author{[AUTHOR NAME]}
\date{October 9, 2026}
\begin{document}
\begin{abstract}
OpenAI's preprint \emph{A finite Smith--Toda complex $V(4)$ at the prime $1009$} chooses
$p=1009$ to make its finite inequalities generous. We check that every finite claim in that
proof holds, as written with $1009$ replaced by $p$, for all primes $p\ge19$. At
$p=11,13,17$ three inequalities fail, but the statements they were used for can be verified
directly by exact enumeration of May $E_1$ pages. \textbf{Conditional on the correctness of the
original $p=1009$ manuscript, which we did not re-verify}, $V(4)$ therefore exists at every prime
$p\ge11$. The argument is OpenAI's; we only re-select the prime and replace three
inequalities by exact enumeration, and we list every changed line. Developed with AI assistance
(Anthropic Claude, Opus 5.5) under [AUTHOR NAME]'s direction; independent adversarial reviews
were run by separate AI sessions and are not human peer review.
\end{abstract}
\maketitle

\section{Statement}
Let \cite{OAI308} denote
\path{preprints/A-finite-Smith-Toda-complex-V4-at-the-prime-1009-September-23-2026}
of \url{https://github.com/openai/math} (commit \texttt{08fcd61}; line numbers refer to
\texttt{build/sections/*.tex}). Its main theorem: at $p=1009$ there is a finite $p$-local spectrum
$X$ with $BP_*X\cong BP_*/(p,v_1,v_2,v_3,v_4)$ as $BP_*BP$-comodules.

\begin{theorem}\label{thm}
Assume that the proof in \cite{OAI308} is correct at $p=1009$. Then for every prime $p\ge 11$ there is
a finite $p$-local spectrum $X$ with $BP_*X\cong BP_*/(p,v_1,v_2,v_3,v_4)$ as graded
$BP_*BP$-comodules, with the canonical quotient coaction and generator in degree zero.
\end{theorem}

\section{Where $p$ enters}
The prime enters \cite{OAI308} only through finite statements about the May $E_1$ page
$\Lambda(h_{i,j})\otimes\mathbb F_p[b_{i,j}]$, with $h_{i,j}$ in bidegree $(1,p^ju_i)$,
$b_{i,j}$ in $(2,p^{j+1}u_i)$ and $u_i=(p^i-1)/(p-1)$, together with explicit polynomial and
divisibility inequalities used to prove those statements. We extracted every such inequality
(33 in total; \texttt{ineq308.py} lists them with line numbers) and determined, for each,
the primes at which it holds. All hold for $p\ge19$. Exactly three fail somewhere in
$\{11,13,17\}$:
\begin{center}
\begin{tabular}{lll}
\toprule
line & inequality & holds for\\\midrule
\texttt{02} L109--113 & $p^2-\tfrac{27}{2}p-9>0$ & $p\ge17$\\
\texttt{02} L154--155 & no nonzero multiple of $p$ in $[-14,3]$ & $p\ge17$\\
\texttt{03} L161--163 & $2p-19>p-2$ & $p\ge19$\\\bottomrule
\end{tabular}
\end{center}
The first two are used only to prove Lemma \texttt{low:algebraic-vanishing}: certain cobar groups
vanish because their May $E_1$ page is empty. The third is used only in the exhaustion proof of
Table \texttt{may:n5-table}. At $p=11,13,17$ we verify these conclusions directly. We enumerate every
$E_1$ monomial in every relevant bidegree: $(qk+1,u_n+k)$ for $n\le3$, and $(qk+1-\rho,U+k)$ for every
profile $e\in\{0,1,2\}^{n+1}$ with $\rho\ge2$, for all $k$ up to the normalized cutoff of
\texttt{low:normalized-cutoff}; also the bidegrees $(q+1,u_4+1)$ and $(q+1,u_5+1)$. The vanishing groups
have no monomials. The table at $(q+1,u_5+1)$ has exactly the five entries of \cite{OAI308}, with
exponents $p-2,p-3,p-3,p-3,p-4$ in place of $1007,\dots,1005$, and $d_1$ is injective over $\mathbb F_p$
on the exterior-length-$3$ and length-$5$ spans. In the same runs we re-verify the other finite
tables the proof uses at each $p$:
\begin{itemize}
\item the $v_4$-target lists, non-empty exactly at $k=1,2,p-1,p$, with the paper's monomials;
\item the Frobenius quotient-group table;
\item the small-stem lemma, with $d_1$-cohomology $0$ at $(6,W+1)$;
\item the page bound $qp+1-S=2p-6>0$, $N-K=p-3>0$.
\end{itemize}
All other uses of $1009$ in \cite{OAI308} are numerical instantiations of these symbolic
inequalities (the README lists them line by line). The adversarial review read every section,
including the appendix, and found no further dependence on the size of $p$, in particular no
argument of the form ``$p$ is large, hence a group vanishes''.

\begin{remark}[Controls]
At $p=7$ the same enumeration finds non-empty product-obstruction pages, for example
$y\,A^{16}$ in bidegree $(33,119)$ for the profile $e=(0,2,0,2)$, $k=3$. So the mechanism fails at
$p=7$. This is consistent with Toda's requirement $p\ge11$ for a unital product on $V(3)$ and with Nave's
theorem that $V((p+1)/2)=V(4)$ does not exist at $p=7$. At $p=1009$ the enumeration reproduces the
paper's tables verbatim (positive control).
\end{remark}

\begin{remark}[What is not claimed]
We do not verify the homotopy-theoretic content of \cite{OAI308}: the adjusted-filtration injectivity
(\S3), the Frobenius left-hand side (\S4), the extension of powers (\S5) and the appendix pairing.
Theorem~\ref{thm} is exactly as reliable as that proof. Nothing is claimed for $p=3,5$.
\end{remark}

\section{Certificates}
\texttt{may308.py} (direct recursive enumeration, exact integers, $d_1$ ranks over $\mathbb F_p$) and
\texttt{may308\_dp.py} (knapsack DP counts) form the primary computation. \texttt{independent/sweep.py} and
\texttt{independent/d1check.py} were written separately; they use z3 exact enumeration with blocking clauses
and an int64 DP modulo $2^{61}-1$. All agree at $p=7,11,13,17$ and at $p=1009$ (sampled). \texttt{may308.py}
also passes at $p=19,\dots,47$.

\begin{thebibliography}{9}
\bibitem{OAI308} OpenAI, \emph{A finite Smith--Toda complex $V(4)$ at the prime $1009$}, OpenAI Math
Release preprint, September 23, 2026,
\url{https://github.com/openai/math/tree/main/preprints/A-finite-Smith-Toda-complex-V4-at-the-prime-1009-September-23-2026}.
\bibitem{Nave} L.~S.~Nave, \emph{The Smith--Toda complex $V((p+1)/2)$ does not exist}, Ann.\ of Math.\ 171 (2010), 491--509.
\bibitem{Toda} H.~Toda, \emph{On spectra realizing exterior parts of the Steenrod algebra}, Topology 10 (1971), 53--65.
\end{thebibliography}
\end{document}
